\documentclass[11pt,a4paper]{article}
% Shared preamble for the qi26_25 weekly deliverable reports.
%
% Each weekly report is a short standalone document that inputs this file, so
% that formatting stays consistent across all six without duplicating the
% package list in every one. Change something here and every report follows.
%
% Usage in a weekly report:
%     \documentclass[11pt,a4paper]{article}
%     % Shared preamble for the qi26_25 weekly deliverable reports.
%
% Each weekly report is a short standalone document that inputs this file, so
% that formatting stays consistent across all six without duplicating the
% package list in every one. Change something here and every report follows.
%
% Usage in a weekly report:
%     \documentclass[11pt,a4paper]{article}
%     % Shared preamble for the qi26_25 weekly deliverable reports.
%
% Each weekly report is a short standalone document that inputs this file, so
% that formatting stays consistent across all six without duplicating the
% package list in every one. Change something here and every report follows.
%
% Usage in a weekly report:
%     \documentclass[11pt,a4paper]{article}
%     \input{weekly_preamble}
%     \weekhead{1}{Vision Foundation and Data Encoding}{Week of ...}
%     ... body ...
%     \end{document}

\usepackage[utf8]{inputenc}
\usepackage[T1]{fontenc}
\usepackage[margin=2.5cm]{geometry}
\usepackage{amsmath,amssymb}
\usepackage{booktabs}
\usepackage{graphicx}
\usepackage{enumitem}
\usepackage{listings}
\usepackage{xcolor}
\usepackage[hidelinks]{hyperref}
\usepackage{titlesec}
\usepackage{fancyhdr}

\definecolor{codebg}{RGB}{247,247,250}
\definecolor{codekw}{RGB}{20,90,160}
\definecolor{codecm}{RGB}{110,120,130}
\definecolor{accent}{RGB}{20,90,160}

\lstset{
  backgroundcolor=\color{codebg},
  basicstyle=\ttfamily\footnotesize,
  keywordstyle=\color{codekw}\bfseries,
  commentstyle=\color{codecm}\itshape,
  breaklines=true,
  frame=single,
  rulecolor=\color{gray!30},
  language=Python,
  showstringspaces=false,
  xleftmargin=6pt,
  framexleftmargin=6pt
}

\setlist{itemsep=2pt,topsep=4pt}
\titleformat{\section}{\normalfont\large\bfseries\color{accent}}{\thesection}{0.8em}{}
\titleformat{\subsection}{\normalfont\normalsize\bfseries}{\thesubsection}{0.8em}{}

\pagestyle{fancy}
\fancyhf{}
\fancyfoot[C]{\small\thepage}
\fancyfoot[R]{\scriptsize\texttt{github.com/Pushkar0997/qi26\_25}}
\renewcommand{\headrulewidth}{0pt}

% \weekhead{number}{title}{deliverable as stated in the project brief}
\newcommand{\weekhead}[3]{%
  \begin{center}
    {\scriptsize\textsf{QINTERN 2026 \textbullet\ QWORLD \textbullet\ WEEKLY DELIVERABLE}}\\[0.25cm]
    {\LARGE\bfseries Week #1: #2}\\[0.35cm]
    {\normalsize Hybrid QML and Quantum String Algorithms for Document Extraction}\\[0.25cm]
    {\small Pushkar Kumar \quad\textbullet\quad Mentor: Potluri Krishna Priyatham}\\[0.4cm]
    \fbox{\begin{minipage}{0.88\textwidth}
      \small\textbf{Deliverable as specified in the project brief:} #3
    \end{minipage}}
  \end{center}
  \vspace{0.4cm}
}

% A boxed note for caveats and honest limitations.
\newcommand{\honestnote}[1]{%
  \vspace{0.2cm}
  \noindent\fcolorbox{gray!40}{gray!8}{%
    \begin{minipage}{0.97\textwidth}\small #1\end{minipage}}
  \vspace{0.2cm}
}

%     \weekhead{1}{Vision Foundation and Data Encoding}{Week of ...}
%     ... body ...
%     \end{document}

\usepackage[utf8]{inputenc}
\usepackage[T1]{fontenc}
\usepackage[margin=2.5cm]{geometry}
\usepackage{amsmath,amssymb}
\usepackage{booktabs}
\usepackage{graphicx}
\usepackage{enumitem}
\usepackage{listings}
\usepackage{xcolor}
\usepackage[hidelinks]{hyperref}
\usepackage{titlesec}
\usepackage{fancyhdr}

\definecolor{codebg}{RGB}{247,247,250}
\definecolor{codekw}{RGB}{20,90,160}
\definecolor{codecm}{RGB}{110,120,130}
\definecolor{accent}{RGB}{20,90,160}

\lstset{
  backgroundcolor=\color{codebg},
  basicstyle=\ttfamily\footnotesize,
  keywordstyle=\color{codekw}\bfseries,
  commentstyle=\color{codecm}\itshape,
  breaklines=true,
  frame=single,
  rulecolor=\color{gray!30},
  language=Python,
  showstringspaces=false,
  xleftmargin=6pt,
  framexleftmargin=6pt
}

\setlist{itemsep=2pt,topsep=4pt}
\titleformat{\section}{\normalfont\large\bfseries\color{accent}}{\thesection}{0.8em}{}
\titleformat{\subsection}{\normalfont\normalsize\bfseries}{\thesubsection}{0.8em}{}

\pagestyle{fancy}
\fancyhf{}
\fancyfoot[C]{\small\thepage}
\fancyfoot[R]{\scriptsize\texttt{github.com/Pushkar0997/qi26\_25}}
\renewcommand{\headrulewidth}{0pt}

% \weekhead{number}{title}{deliverable as stated in the project brief}
\newcommand{\weekhead}[3]{%
  \begin{center}
    {\scriptsize\textsf{QINTERN 2026 \textbullet\ QWORLD \textbullet\ WEEKLY DELIVERABLE}}\\[0.25cm]
    {\LARGE\bfseries Week #1: #2}\\[0.35cm]
    {\normalsize Hybrid QML and Quantum String Algorithms for Document Extraction}\\[0.25cm]
    {\small Pushkar Kumar \quad\textbullet\quad Mentor: Potluri Krishna Priyatham}\\[0.4cm]
    \fbox{\begin{minipage}{0.88\textwidth}
      \small\textbf{Deliverable as specified in the project brief:} #3
    \end{minipage}}
  \end{center}
  \vspace{0.4cm}
}

% A boxed note for caveats and honest limitations.
\newcommand{\honestnote}[1]{%
  \vspace{0.2cm}
  \noindent\fcolorbox{gray!40}{gray!8}{%
    \begin{minipage}{0.97\textwidth}\small #1\end{minipage}}
  \vspace{0.2cm}
}

%     \weekhead{1}{Vision Foundation and Data Encoding}{Week of ...}
%     ... body ...
%     \end{document}

\usepackage[utf8]{inputenc}
\usepackage[T1]{fontenc}
\usepackage[margin=2.5cm]{geometry}
\usepackage{amsmath,amssymb}
\usepackage{booktabs}
\usepackage{graphicx}
\usepackage{enumitem}
\usepackage{listings}
\usepackage{xcolor}
\usepackage[hidelinks]{hyperref}
\usepackage{titlesec}
\usepackage{fancyhdr}

\definecolor{codebg}{RGB}{247,247,250}
\definecolor{codekw}{RGB}{20,90,160}
\definecolor{codecm}{RGB}{110,120,130}
\definecolor{accent}{RGB}{20,90,160}

\lstset{
  backgroundcolor=\color{codebg},
  basicstyle=\ttfamily\footnotesize,
  keywordstyle=\color{codekw}\bfseries,
  commentstyle=\color{codecm}\itshape,
  breaklines=true,
  frame=single,
  rulecolor=\color{gray!30},
  language=Python,
  showstringspaces=false,
  xleftmargin=6pt,
  framexleftmargin=6pt
}

\setlist{itemsep=2pt,topsep=4pt}
\titleformat{\section}{\normalfont\large\bfseries\color{accent}}{\thesection}{0.8em}{}
\titleformat{\subsection}{\normalfont\normalsize\bfseries}{\thesubsection}{0.8em}{}

\pagestyle{fancy}
\fancyhf{}
\fancyfoot[C]{\small\thepage}
\fancyfoot[R]{\scriptsize\texttt{github.com/Pushkar0997/qi26\_25}}
\renewcommand{\headrulewidth}{0pt}

% \weekhead{number}{title}{deliverable as stated in the project brief}
\newcommand{\weekhead}[3]{%
  \begin{center}
    {\scriptsize\textsf{QINTERN 2026 \textbullet\ QWORLD \textbullet\ WEEKLY DELIVERABLE}}\\[0.25cm]
    {\LARGE\bfseries Week #1: #2}\\[0.35cm]
    {\normalsize Hybrid QML and Quantum String Algorithms for Document Extraction}\\[0.25cm]
    {\small Pushkar Kumar \quad\textbullet\quad Mentor: Potluri Krishna Priyatham}\\[0.4cm]
    \fbox{\begin{minipage}{0.88\textwidth}
      \small\textbf{Deliverable as specified in the project brief:} #3
    \end{minipage}}
  \end{center}
  \vspace{0.4cm}
}

% A boxed note for caveats and honest limitations.
\newcommand{\honestnote}[1]{%
  \vspace{0.2cm}
  \noindent\fcolorbox{gray!40}{gray!8}{%
    \begin{minipage}{0.97\textwidth}\small #1\end{minipage}}
  \vspace{0.2cm}
}


\begin{document}
\weekhead{3}{Integration and Raw Extraction}
{Working Python script extracting text from sample documents.}

\section{Objective}

Connect the quantum vision output to a digital text stream: interface the
quantum-processed features with a lightweight classical OCR backend, and produce
a script that takes a document image and returns extracted text.

\section{Pipeline}

Nine stages, of which two are quantum:

\begin{center}
\small
\begin{tabular}{@{}rll@{}}
\toprule
\# & Stage & Description \\
\midrule
1 & preprocess & median filter, Otsu threshold, speckle removal \\
2 & segment & projection profiles $\rightarrow$ character boxes \\
3 & crop & pad to square, resize to $8\times8$, normalise \\
4--5 & \textbf{quantum} & \textbf{$R_Y$ encoding + quanvolutional layer $\rightarrow$ 160 features} \\
6 & classify & multinomial logistic regression $\rightarrow$ characters \\
7 & assemble & rebuild lines, insert spaces \\
8 & locate & identify the document identifier field \\
9 & \textbf{quantum} & \textbf{Grover comparator search (Week 4)} \\
\bottomrule
\end{tabular}
\end{center}

The classifier at stage 6 is the ``lightweight classical OCR back-end'' the brief
specifies: feature extraction is quantum, the decision layer is not.

\section{Classical Front End}

Three preprocessing operations, each addressing a specific failure.

\textbf{Median filter} removes isolated speckle while preserving edges, unlike a
Gaussian blur which would smear thin strokes.

\textbf{Otsu threshold} selects the intensity cut maximising between-class
variance. A fixed threshold was rejected because the five tiers differ greatly in
overall brightness and any constant fails on at least one.

\textbf{Speckle removal} drops connected components too small to be a stroke.
This is not cosmetic: without it the noisy tiers place sub-threshold pixels in
\emph{every} row of the page, so line detection found a single 420-row ``line''
and segmentation collapsed entirely.

Segmentation then sums ink along rows to find text lines, and within each line
along columns to find characters. The row threshold is adaptive---a real text
line marks a meaningful fraction of the page width---because a fixed
``more than one pixel'' test allowed residual noise to merge the page into one
line on degraded scans.

\subsection{Two Fixes Worth Recording}

\textbf{Aspect ratio.} Each character box is padded to square \emph{before}
resizing to $8\times8$. Resizing directly destroys aspect ratio, and for OCR this
is not cosmetic: a colon is two small dots in a narrow tall box, and stretched to
a square it becomes two thick horizontal bands closely resembling an
\texttt{E}. Before this was fixed, \texttt{C}$\rightarrow$\texttt{:} and
\texttt{E}$\rightarrow$\texttt{:} confusions dominated the error and CER was
32\%.

\textbf{Merged glyph splitting.} On low-DPI scans, downsampling plus blur closes
the gap between adjacent glyphs so the column profile sees no separator. The
obvious fix---cutting a too-wide box into equal slices---was tried and was
\emph{wrong}: on degraded tiers it over-segmented badly, cutting single
characters apart and pushing error above 100\% through pure insertions. The
working version cuts only at genuine valleys in the ink profile and leaves a box
intact if it has none, so a legitimately wide \texttt{M} survives.

\section{Results}

\begin{table}[h]
\centering
\caption{End-to-end extraction, all 12 documents per tier.}
\begin{tabular}{@{}lrrr@{}}
\toprule
Tier & CER & Segmentation ratio & ID recovered exactly \\
\midrule
clean\_digital & 6.6\% & 0.95 & \textbf{100\%} \\
clean\_scan & 8.5\% & 0.93 & 58\% \\
noisy\_scan & 79.5\% & 0.34 & 0\% \\
clear\_handwriting & 51.8\% & 0.70 & 0\% \\
degraded\_handwriting & 85.7\% & 0.99 & 0\% \\
\bottomrule
\end{tabular}
\end{table}

Two metrics are reported because they diverge. Character error rate measures
average character quality; identifier recovery measures whether the pipeline did
its actual job, for which one wrong character makes the result useless.
\texttt{clean\_scan} reads at 8.5\% CER---visually almost clean---yet yields a
usable identifier in only 58\% of documents.

\honestnote{%
\textbf{Three of five tiers fail outright.} Noisy scans and both handwriting
tiers have character error rates of 52--86\% and recover no identifiers at all.
The pipeline is usable on clean digital input and not usable on degraded input,
and the tier-by-tier breakdown is reported rather than an average that would
obscure it.}

\section{The Largest Error Source Was a Data Bug}

The biggest single improvement in the entire project came from a data-handling
defect, not from any modelling change.

\textbf{How it surfaced.} Switching the feature stage from the 64-dimensional
marginal readout to the 160-dimensional $+ZZ$ readout \emph{raised} isolated
character accuracy but made end-to-end CER \emph{worse}. A strictly better
feature set producing strictly worse output is impossible if both are measured on
the same input distribution---so they were not.

\textbf{The cause.} Training crops were cut from exact glyph bounds recorded
during dataset generation; inference crops come from projection segmentation.
These are measurably different distributions: on \texttt{clean\_scan},
ground-truth boxes have median size $16\times19$ px against segmented boxes at
$12\times15$ px. The higher-capacity features fitted the training distribution
more tightly and transferred less well.

\textbf{The fix.} Training crops are now built by running the real segmentation
path over training pages and labelling each detected box from ground truth by
overlap, with ambiguous detections dropped rather than guessed.

\begin{table}[h]
\centering
\caption{Effect of correcting the train/serve skew.}
\begin{tabular}{@{}lrrr@{}}
\toprule
Tier & CER before & CER after & Reduction \\
\midrule
clean\_digital & 42.9\% & \textbf{6.6\%} & $6.5\times$ \\
clean\_scan & 29.7\% & \textbf{8.5\%} & $3.5\times$ \\
clear\_handwriting & 74.5\% & \textbf{51.8\%} & $1.4\times$ \\
\bottomrule
\end{tabular}
\end{table}

\subsection{A Caution About the Accuracy Figure}

Training on segmented crops \emph{lowers} reported held-out character accuracy
from 92.6\% to 67.7\%. This is not a regression. The two numbers are measured on
different crop distributions, and the lower one is measured on the harder,
realistic distribution.

The comparison to draw: a model scoring 92.6\% on clean bounds produced 42.9\%
document CER, while a model scoring 67.7\% on realistic bounds produced 6.6\%.

\honestnote{%
\textbf{Isolated-character accuracy was an actively misleading proxy for
end-to-end accuracy, and optimising it drove the pipeline in the wrong
direction.} Any comparison of feature extractors on pre-cut crops---including
the Week 2 comparison---measures something narrower than pipeline quality. This
is the most transferable finding in the project.}

\section{Where the Remaining Error Lies}

On clean tiers, segmentation recovers 93--95\% of characters and isolated
character accuracy approaches 99\%, yet CER is 6.6--8.5\%. The residual is
therefore not classification but missed or merged glyphs and word-spacing
reconstruction.

\textbf{The classical front end, not either quantum stage, is the accuracy
bottleneck on clean documents.}

\section{Deliverable Status}

Complete. A working script takes a document image and produces extracted text,
run successfully across all five tiers.

\vspace{0.3cm}
\noindent\textbf{Artifacts:} \texttt{integration/pipeline.py},
\texttt{notebooks/04\_full\_pipeline.ipynb}

\vspace{0.2cm}
\noindent\textbf{Usage:}
\texttt{python integration/pipeline.py --doc data/processed/clean\_scan/doc\_000.png --pattern 38}

\end{document}
